\begin{remark}\label{rem:newEss}
With notation as in Lemma \ref{lem:SchubertDel}, it is possible that $\Ess(w') \not\subseteq \Ess(w)$, as is the case if $w = 215634$ and $(i,j) = (4,4)$.  In that case, $(3,4), (4,3) \in \Ess(w') \setminus \Ess(w)$.  

In general, the only possible elements of $\Ess(w') \setminus \Ess(w)$ are $(i-1,j)$ and $(i,j-1)$.  Indeed, suppose $(a,b) \in \Ess(w') \setminus \Ess(w)$. By the definition of $\Ess(w')$, $(a,b) \in D_{w'}$ and $(a+1,b), (a,b+1) \notin D_{w'}$.  Because $D_{w'} \subset D_w$, the assumption $(a,b) \notin \Ess(w)$ implies $(a+1,b) \in \Ess(w)$ or $(a,b+1) \in \Ess(w)$.  The fact that $D_w \setminus D_{w'} = \{(i,j)\}$ then implies that $(a+1,b) = (i,j)$ or $(a,b+1) = (i,j)$.
\end{remark}

\begin{corollary}\label{cor:GrobnerDel}
If $w$ has no obstruction of Type $1$, $2$, or $3$, $(i,j)$ is a lower outside corner of $D_w$ corresponding to the variable $y = z_{i,j}$, and $I_{w'} =N_{y,I_w}$, then $w'$ has no obstruction of Type $1$, $2$, or $3$.  
\end{corollary}

\begin{proof}
If $(i,j) \in \Dom(w)$, then, because there are no diagram boxes southeast of $(i,j)$ by assumption, $(i,j)$ cannot be one of the boxes involved in any of the obstructions.  In that case, any set of diagram boxes realizing an obstruction in $w'$ would also realize an obstruction in $w$.  Hence we assume that $(i,j) \notin \Dom(w)$, in which case $\Dom(w) = \Dom(w')$.

Suppose that $w'$ has an obstruction of Type $1$ with $(r,s) \in \Dom(w') \cap \Ess(w')$ and $(a,b)$, $(a'b') \in D_{w'}$ strictly southeast of $(r,s)$ with $a' \neq a$ and $b' \neq b$.  Because $\Dom(w') = \Dom(w)$ and $D_{w'} \subset D_w$, the diagram boxes $(r,s)$, $(a,b)$, $(a',b')$ also constitute an obstruction of Type $1$ in $w$ as well.

Next suppose that $w'$ has an obstruction of Type $2$.  Assume without loss of generality that the Type $2$ obstruction has the form  $(r,s) \in \Dom(w') \cap \Ess(w')$ and $(a,b), (a,b') \in \Ess(w')$ strictly southeast of $(r,s)$ with $\max_k \{(k,b) \in \Dom(w')\} = \max_k \{(k,b') \in \Dom(w') \}$ and $b<b'$.  In this case neither $(a,b)$ nor $(r,s)$ may be a lower outside corner of $w'$.   In particular, neither $(r,s)$ nor $(a,b)$ is equal to $(i,j)$, and so $(r,s) \in \Dom(w) \cap \Ess(w)$ and $(a,b) \in \Ess(w)$.  If $(a,b') \in \Ess(w)$ also, then $(r,s)$, $(a,b)$, and $(a,b')$ also constitute a Type $2$ obstruction in $w$.  If $(a,b') \notin \Ess(w)$, then either $(a+1,b') = (i,j) \in \Ess_w$, in which case $(r,s)$, $(a,b)$, and $(a+1,b')$ together constitute a Type $1$ obstruction, or $(a,b'+1) = (i,j) \in \Ess(w)$.  

If $(a,b'+1) = (i,j) \in \Ess(w)$, set $m = \max_k \{(k,b) \in \Dom(w')\} = \max_k \{(k,b') \in \Dom(w')\}$.  Because $\Dom(w) = \Dom(w')$, we have $m = \max_k \{(k,b) \in \Dom(w)\}$.  We claim that $m = \max_k \{(k,b'+1) \in \Dom(w)\}$.   Because $b'+1>b'$ and $\Dom(w)$ forms a partition shape, $m \geq \max_k \{(k,b'+1) \in \Dom(w)\}$.  If $m > \max_k \{(k,b'+1) \in \Dom(w)\}$, then $(m,b'+1) \notin \Dom(w)$.  But $(m,b') \in \Dom(w') = \Dom(w)$.  Therefore, if $(m,b'+1) \notin \Dom(w)$, the visualization of $D_w$ must have a $\bullet$ in column $b'+1$ weakly north of row $m$.  Clearly $m<a$ since $(a,b) \notin \Dom(w')$.  But the $\bullet$ in column $b'+1 = j$ must be south of row $a = i$ because $(i,j) \in D_w$. Hence, $m = \max_k \{(k,b'+1) \in \Dom(w)\}$ and $(r,s)$, $(a,b)$, and $(a, b'+1)$ together constitute a Type $3$ obstruction in $w$. 

Finally suppose that $w'$ has an obstruction of Types $3$.  Suppose that $(a,b), (a',b') \in \Ess(w') \setminus \Dom(w')$ with $(a',b')$ strictly southeast of $(a,b)$.  If $(a',b') \in \Ess(w)$, then $(a,b)$ and $(a',b')$ constitute a Type $3$ obstruction in $w$ also.  Otherwise, it must be that $(i,j) = (a'+1,b')$ or $(i,j) = (a',b+1)$.  In either case, $(a,b)$ and $(i,j)$ constitute a Type $3$ obstruction in $w$.
\end{proof}
